\documentclass[11pt,a4paper]{article}
\usepackage{../../intuitmath}
\hypersetup{pdftitle={Action and the Calculus of Variations},pdfauthor={Mathematics and Physics: Intuition First}}
\renewcommand{\sectionmark}[1]{\markboth{#1}{}}

\begin{document}
\tutorialtitle{Action and the Calculus of Variations}{Why mechanics compares whole paths}{From a small change of a curve to the Euler--Lagrange equation}
\tableofcontents
\clearpage

\section{The new kind of question}

Ordinary calculus asks how a number changes when we change another number. If
\[
f(x)=x^2,
\]
then $f'(x)$ measures the response to a small change $x\mapsto x+\varepsilon$.

The calculus of variations asks a similar question about an entire function. A path
\[
q:[t_0,t_1]\longrightarrow\R
\]
is the object being changed. A \emph{functional} assigns one number to the whole path. The central example in mechanics is the action
\[
\boxed{S[q]=\int_{t_0}^{t_1}L(q,\dot q,t)\,\dd t.}
\]

The square brackets are a useful reminder: $S$ consumes a function $q(t)$, not merely its value at one instant.

\begin{intuition}[A landscape whose points are curves]
Imagine that every admissible path is one point in a vast ``space of paths''. The action gives a height to each point. A physical path is stationary when every allowed infinitesimal direction through this space produces no first-order change in height.
\end{intuition}

\section{Why an integral over time?}

The Lagrangian $L(q,\dot q,t)$ measures something locally, at one instant. The action adds that local contribution along the entire history. For an ordinary conservative particle,
\[
L=T-V.
\]
The action therefore balances motion and potential along a path rather than choosing the smallest energy at every instant.

\begin{warningbox}[Stationary does not always mean smallest]
The physical path makes the first variation vanish. It may give a local minimum, a local maximum, or a saddle value of the action. ``Principle of stationary action'' is more accurate than ``principle of least action''.
\end{warningbox}

\section{How to vary a path}

Choose a reference path $q(t)$ and a smooth shape $\eta(t)$. Define a one-parameter family
\[
q_\varepsilon(t)=q(t)+\varepsilon\eta(t).
\]
The small number $\varepsilon$ controls the size of the change; $\eta$ controls its shape.

If the endpoints are fixed, every neighbouring path must begin and end at the same places, so
\[
\boxed{\eta(t_0)=\eta(t_1)=0.}
\]

\begin{center}
\begin{tikzpicture}[x=1.25cm,y=1.05cm,>=Latex]
  \draw[->,ink!65] (0,0) -- (6.5,0) node[right] {$t$};
  \draw[->,ink!65] (0,-0.2) -- (0,3.5) node[above] {$q$};
  \draw[blue,line width=1.6pt] plot[smooth] coordinates {(0.4,.7) (1.4,1.2) (2.5,2.15) (3.7,2.35) (5.0,1.55) (6.0,1.1)};
  \draw[teal,dashed,line width=1.4pt] plot[smooth] coordinates {(0.4,.7) (1.4,1.55) (2.5,1.75) (3.7,2.75) (5.0,1.3) (6.0,1.1)};
  \fill[ink] (0.4,.7) circle (2pt) (6.0,1.1) circle (2pt);
  \node[blue] at (4.7,2.15) {$q(t)$};
  \node[teal] at (3.9,3.02) {$q(t)+\varepsilon\eta(t)$};
  \draw[<->,gold,line width=1pt] (2.5,1.78) -- (2.5,2.12) node[midway,right] {$\varepsilon\eta$};
\end{tikzpicture}
\end{center}

Because differentiation is linear,
\[
\dot q_\varepsilon=\dot q+\varepsilon\dot\eta.
\]
Changing the path therefore changes both the position argument and the velocity argument of $L$.

\section{The first variation}

Insert the varied path into the action:
\[
S[q_\varepsilon]
=\int_{t_0}^{t_1}L(q+\varepsilon\eta,\dot q+\varepsilon\dot\eta,t)\,\dd t.
\]
Now $S[q_\varepsilon]$ is an ordinary function of the number $\varepsilon$. A stationary path satisfies
\[
\left.\td{}{\varepsilon}S[q_\varepsilon]\right|_{\varepsilon=0}=0.
\]

Using the chain rule,
\begin{align*}
\left.\td{}{\varepsilon}S[q_\varepsilon]\right|_{\varepsilon=0}
&=\int_{t_0}^{t_1}
\left(
\pd{L}{q}\eta+\pd{L}{\dot q}\dot\eta
\right)\dd t.
\end{align*}
This expression is called the \emph{first variation} and is written $\delta S$.

\section{The decisive integration by parts}

The second term contains $\dot\eta$, but we want the arbitrary function itself. Integrate by parts:
\begin{align*}
\int_{t_0}^{t_1}\pd{L}{\dot q}\dot\eta\,\dd t
&=\left[\pd{L}{\dot q}\eta\right]_{t_0}^{t_1}
-\int_{t_0}^{t_1}\td{}{t}\left(\pd{L}{\dot q}\right)\eta\,\dd t.
\end{align*}
The boundary term vanishes because $\eta(t_0)=\eta(t_1)=0$. Hence
\[
\delta S=\int_{t_0}^{t_1}
\left[
\pd{L}{q}-\td{}{t}\left(\pd{L}{\dot q}\right)
\right]\eta(t)\,\dd t.
\]

The function $\eta$ may be chosen freely in the interior. The only way this integral can vanish for every such choice is for its coefficient to vanish point by point.

\begin{resultbox}[Euler--Lagrange equation]
\[
\boxed{\td{}{t}\left(\pd{L}{\dot q}\right)-\pd{L}{q}=0.}
\]
One integration by parts converts a statement about every nearby path into a differential equation for the stationary path.
\end{resultbox}

\section{Example: the free particle}

Let
\[
L=\frac12m\dot q^2.
\]
Then
\[
\pd{L}{\dot q}=m\dot q,
\qquad
\pd{L}{q}=0.
\]
The Euler--Lagrange equation gives
\[
m\ddot q=0.
\]
Thus the stationary path has constant velocity:
\[
q(t)=At+B.
\]

\begin{intuition}[Why a straight world-line?]
With no potential, unnecessary changes of velocity only increase the kinetic contribution. Between fixed endpoints and fixed times, the stationary path is the line of uniform motion.
\end{intuition}

\section{Example: the harmonic oscillator}

For a mass on a spring,
\[
L=\frac12m\dot q^2-\frac12kq^2.
\]
The two required derivatives are
\[
\pd{L}{\dot q}=m\dot q,
\qquad
\pd{L}{q}=-kq.
\]
Therefore
\[
\boxed{m\ddot q+kq=0.}
\]
The familiar restoring-force equation has emerged from stationarity of a whole-path quantity.

\section{Many coordinates}

If the configuration is described by $q_1,\ldots,q_n$, vary each independently:
\[
q_i\mapsto q_i+\varepsilon\eta_i.
\]
The result is one equation per coordinate:
\[
\boxed{\td{}{t}\left(\pd{L}{\dot q_i}\right)-\pd{L}{q_i}=0,
\qquad i=1,\ldots,n.}
\]
This is why a single scalar function $L$ can encode the dynamics of a many-coordinate system.

\section{What if the endpoint is free?}

When an endpoint value is not fixed, the corresponding value of $\eta$ need not vanish. The boundary term
\[
\left[\pd{L}{\dot q}\eta\right]_{t_0}^{t_1}
\]
must then be handled explicitly. For example, if $q(t_1)$ is free while $t_1$ is fixed, stationarity requires
\[
\left.\pd{L}{\dot q}\right|_{t_1}=0.
\]
Such conditions are called \emph{natural boundary conditions}. They are not extra guesses; they are the part of the variational calculation that fixed endpoints previously hid.

\section{Adding a total derivative}

Suppose
\[
\widetilde L=L+\td{}{t}F(q,t).
\]
Then
\[
\widetilde S=S+F(q(t_1),t_1)-F(q(t_0),t_0).
\]
For fixed endpoints the added quantity is fixed, so it does not change the stationary paths. Different-looking Lagrangians can therefore encode the same equations of motion.

\section{From particles to fields}

For a field $\phi(x,t)$, the action becomes an integral over space and time:
\[
S[\phi]=\int \mathcal L(\phi,\partial_\mu\phi,x)\,\dd^{d+1}x.
\]
Varying the field and integrating by parts in every coordinate gives
\[
\boxed{
\pd{\mathcal L}{\phi}
-\partial_\mu\left(\pd{\mathcal L}{(\partial_\mu\phi)}\right)=0.
}
\]
The logical structure is unchanged: vary the whole object, move derivatives off the arbitrary variation, and require the remaining coefficient to vanish.

\section{First variation versus second variation}

The condition $\delta S=0$ identifies a stationary path. To classify it, examine the quadratic term in
\[
S[q+\varepsilon\eta]
=S[q]+\varepsilon\,\delta S
+\frac12\varepsilon^2\,\delta^2S+\cdots.
\]
A positive second variation for every nonzero admissible $\eta$ indicates a local minimum. An indefinite second variation indicates a saddle. Mechanics usually begins with the first variation because it already supplies the equations of motion.

\section{A reliable derivation checklist}

\begin{enumerate}
\item State the functional and the allowed class of paths.
\item Write $q_\varepsilon=q+\varepsilon\eta$.
\item Translate endpoint restrictions into conditions on $\eta$.
\item Differentiate with respect to $\varepsilon$ at $\varepsilon=0$.
\item Integrate by parts until derivatives no longer act on $\eta$.
\item Keep or justify every boundary term.
\item Use the arbitrariness of the interior variation.
\item Check the resulting equation against a familiar special case.
\end{enumerate}

\begin{checkBox}[One-sentence summary]
The Euler--Lagrange equation is the pointwise condition forced on a path when the first-order change of an integral functional vanishes for every allowed infinitesimal deformation of that path.
\end{checkBox}

\end{document}
